\chapter*{Resumen}
\phantomsection
\addcontentsline{toc}{chapter}{Resumen}
\markboth{Resumen}{Resumen}

\noindent\textbf{Título original:} \TituloTFG

\vspace{0.4cm}

Harmony Hub es una aplicación móvil orientada a transformar el estado
emocional y contextual del usuario en una experiencia musical más útil para su
bienestar cotidiano. La solución combina un check-in guiado, una lógica de
recomendación contextual, un catálogo musical estructurado, la materialización
de playlists reales en Spotify y un ciclo de aprendizaje apoyado en feedback.
Desde el punto de vista técnico, el sistema se construye sobre un cliente móvil
en Flutter, un backend en FastAPI, persistencia documental en Firestore y un
modelo supervisado a nivel de sesión cuya activación solo se permite cuando
existe evidencia objetiva suficiente.

El proceso de recomendación no depende exclusivamente de un modelo de
aprendizaje automático. En su lugar, el sistema adopta un enfoque híbrido: una
capa heurística garantiza operatividad desde la primera sesión, mientras que un
modelo de regresión logística puede reordenar los modos candidatos cuando se
alcanzan mínimos de datos, cobertura por clase y calidad global. El ranking
musical principal se realiza sobre el catálogo propio \texttt{msd\_tracks},
mientras que Spotify se utiliza principalmente para autenticación, lectura de
contexto musical básico y materialización final de la playlist. Esta decisión
permite mantener utilidad incluso bajo restricciones externas, como cambios en
endpoints, limitaciones de cuota o ausencia de datos de entrenamiento
suficientes.

El proyecto explora así cómo combinar software móvil, APIs externas, datos
contextuales y aprendizaje progresivo para acompañar momentos cotidianos de
regulación emocional sin plantearse como una herramienta clínica. En el estado
actual del entregable, el aprendizaje individual del perfil ya es visible en la
experiencia y deja evidencia longitudinal sobre los cinco estados emocionales
considerados. Además, el modelo supervisado de sesión se encuentra activo bajo
puertas explícitas de calidad y puede aportar reordenación probabilística sin
perder la capacidad de abstención cuando la evidencia es débil. La validación
realizada muestra también que la medición del entorno acústico puede integrarse
en el contexto de recomendación y en la generación de playlists. El resultado
es un prototipo funcional que demuestra la viabilidad de una recomendación
musical explicable, orientada a sesiones y capaz de crear playlists reales en
Spotify con trazabilidad contextual y aprendizaje progresivo.
